\clearpage
\heading{What a complete formal check establishes}\label{sec:14formal}
The comparison starts with an independently readable interface. \texttt{Statement.lean} imports mathlib and defines a circuit by a finite initial segment of a list of formal polynomials. In any product gate the two earlier indices may coincide. The output can occur anywhere in the segment. The threshold is quantified in the order
\[
 \exists c>0\quad \forall p\quad\exists m_0\quad
 \forall F\text{ of characteristic }p\quad\forall m\ge m_0\quad\forall s.
\]
The characteristic-two exception occurs in the permanent implication. A dimension-dependent family of circuits and all field constants are allowed. Thus checking the five public declarations includes the all-circuit quantifier. An implementation restricted to formulas or monotone circuits would need a different interface.

\subhead{Five declarations in two isolated environments}
The challenge lists the goal types and uses a placeholder for each proof. The solution imports the common definitions and a proof module. The challenge does not occur in its transitive imports. Both environments consequently contain the same public names. Loading them together would defeat this separation. The local test prints each type in its own environment with implicit arguments and universes exposed, then compares the two transcripts. For example, \texttt{Solution.separation} has as its target the strict inclusion $\forall F,\ \mathrm{VP}(F)\subsetneq\mathrm{VNP}(F)$. The Hamiltonian-cycle witness and the contradiction at large $m$ are premises established inside its proof.
\begin{center}
\begin{tabular}{@{}ll@{}}
\textbf{Solution declaration} & \textbf{Encoded conclusion}\\
\texttt{mainTheorem} & exact HC and eligible permanent size bounds\\
\texttt{separation} & strict inclusion over every field\\
\texttt{mainTheoremApprox} & Laurent-series approximation bounds\\
\texttt{separationClosure} & VNP escapes the defined closure of VP\\
\texttt{permanent\_over\_real} & exact permanent bound over $\mathbb R$
\end{tabular}
\end{center}
The approximate interface requires every negative Laurent coefficient to vanish and each constant coefficient to equal its target coefficient, monomial by monomial. Machine verification covers these two approximate declarations. The detailed human derivations in this edition concentrate on exact circuits.

\subhead{A boundary exposed by the compiler audit}
The K.3 proof can meet its type even though a stronger sentence about a polynomial on malformed inputs fails. Conversely, finite examples that reproduce each sign in a small tower cannot certify all parameter values. Relating the public definitions, the internal types and the printed arguments therefore remains a separate task even when the entire proof term has been checked.

\smallref{\pub{Statement.lean}{20}{Circuit model} · \pub{Statement.lean}{41}{quantifier order} · \pub{Statement.lean}{50}{approximation} · \pub{Statement.lean}{75}{complexity classes} · \pub{Challenge.lean}{17}{challenge declarations} · \pub{Solution.lean}{25}{solution wrappers}.}

\clearpage
\heading{Two independent executions of the publisher's proof}
\subhead{Rebuilding from source on the review machine}
The source lock fixes the 333 Git blobs of the publisher's project. The local runner computes the complete dependency closure of the four default targets, compares it with the entire set of Lean sources including the axiom-reporting script, and verifies the revisions of the external packages. It starts with an empty project object directory. It keeps a per-module SHA-256 digest and compiler exit code, and refuses a proof-placeholder warning outside the isolated challenge. External library objects come from pinned dependency checkouts. Rebuilding those libraries from their own sources is a different scope.

All 324 sources compiled with Lean 4.33.1. The five signature pairs matched and 22 axiom reports contained only the permitted names. The evidence is recorded in \texttt{pass14\_release.json}. Alongside all module results it holds the signature transcripts and the axiom report for each public result and selected intermediate declarations. These allow a reader to detect a changed source module, a different endpoint or an additional axiom. For example, a solution importing its challenge could otherwise acquire a conclusion through the placeholder axiom.

\subhead{Replaying with a second kernel on Linux}
The review repository's hosted job checks the Git tree and the publisher's verifier-script blob, then builds the default targets in a disposable copy. It invokes the publisher's Comparator driver, which builds the export and verification tools at its specified revisions, confines the subprocesses with landrun, and forces the independent nanoda kernel on. Comparator compares the exported proofs against the trusted challenges and enforces the allowed axiom list. Its driver requires both a successful exit and a complete acceptance marker. The independent job completed successfully. Its log records the acceptance marker for the five-name configuration. The separate receipt identifies the job, exported proof check, tool revisions and artifact.

\subhead{What it means for the original question}
A complete replay concerns all sizes, all field constants and all permitted reuses in the encoded statement. It checks the publisher's proof objects in a fresh execution. The shared formal statement, encoding of arithmetic circuits and standard foundational axioms remain assumptions of interpretation. The stronger independent E.10 coefficient in 14A has its own Lean check and sits outside the five-declaration Comparator configuration. The walkthrough beginning on p.~\pageref{sec:14chain} follows the arbitrary circuit to the two terminal contradictions. Pass 15 checks the diagonal's instruction budgets (p.~\pageref{sec:15k6}). A human reconstruction of the approximate-circuit argument remains open.

\smallref{\newreview{checks/pass14_release.json}{Local module and axiom receipt} · \newreview{checks/pass14_comparator.json}{hosted verification receipt} · \href{https://github.com/nousentllc/VP-VNP-Review/actions/runs/45691617731}{review-run Linux job} · \newreview{passes/14d-formal.md}{scope and reproduction instructions}.}
